\section{The model}\label{sec:model}

\subsection{The decision interval}\label{sec:interval}

The model applies to a listed target whose sale opportunity is publicly understood while its stock still
trades. In this interval, a further bidder can still come in. A disclosed approach, an announced strategic
review, or an open contest can create such an interval. The opportunity must become public before the buyer
set is fixed. Only then can a prospective buyer watch the price during the period that matters. Whether a
given transaction fits is a question about its chronology, and Section~\ref{sec:empirical} states how to date it.

Within the interval, one buyer is already prepared. It may be the party that approached the target or a
bidder that has completed diligence. Its strength is what the public can infer about its acquisition value from the disclosed
record. Another buyer has not decided whether to prepare. Preparation is required for an executable proposal.

Before a buyer can submit a binding offer, it must verify the target's contracts, liabilities, and technology.
It must also arrange financing and approvals and commit evaluation capacity that has other uses. Those costs
are sunk once incurred, and they differ across buyers. Prior familiarity, capacity already in place, or a small deal
relative to the buyer's size make preparation cheap for some buyers and expensive for others. A buyer that
does not prepare cannot bid its current expected value instead; it stays out. The challenger's cost is an
entry cost, not the price of an optional report.

The investor and the buyer know different things. A buyer knows its own integration plans. A specialist
investor may know what the target's customers do, how its technology compares with alternatives, or how
its product demand will change. These are different facts about the same acquisition match, so a more
accurate signal about one can leave the other unknown.

The model gives the investor exact knowledge of the challenger's quality class, though not of its value, and
suppresses the buyer's initial information. Section~\ref{sec:signals} removes both extremes, and
Section~\ref{sec:research} makes the investor's information the result of costly research. The identity of the
further buyer is common knowledge; what the market does not know is how well that buyer fits the target.

Public incumbent strength is a distribution, not a bid. The record can reveal that a strategic acquirer with
clear synergies is at the table. It can also reveal that the prepared party is a financial buyer with a
leveraged capital structure. In each case it need not reveal what either buyer would pay. Strategic and
financial bidders differ systematically in their valuations \citep{GorbenkoMalenko2014}. The model varies that
distribution and keeps the incumbent's realized value private until it bids.

Disclosure records document a costly preparation stage. Imprivata's definitive proxy separates an unsolicited
approach, outreach to potential buyers, and indications of interest conditional on further diligence
\citep{Imprivata2016}. Takeover competition develops before public bids
\citep{BooneMulherin2007,GentryStroup2019}.

\subsection{Values, preparation, and the sale}\label{sec:values}

The target's known standalone value is normalized to zero. Acquisition values and preparation costs are
measured per target share. A common standalone component added to every ownership outcome shifts prices and
payoffs by a constant and leaves incentives unchanged.

All strategic agents are risk neutral. The target share is the only traded claim with information about the
acquisition match, and neither bidder trades it. The investor has no control rights and cannot acquire the
target. The sale mechanism binds all target shares, so shareholder tendering and holdout stay outside the
model. The incumbent's value distribution is public when the sale opportunity becomes visible. Its realized
value stays undisclosed to the challenger before preparation.

Two potential acquirers face the target. The incumbent has already prepared, incurs no further entry cost, and
knows its value before it bids. Its acquisition value \(R\) has a continuous distribution \(F\) on
\([0,\bar r]\), conditional on public information when the sale process begins. The distribution \(F\) is
extended by unity above \(\bar r\). A first-order
stochastic shift of \(F\) strengthens the incumbent. In the main case \(R\) is uniform on \([0,r]\), and an
increase in \(r\) strengthens the incumbent.

The challenger belongs to one of two quality classes, \(j\in\{H,L\}\), with prior \(\pi=\Pr(j=H)\in(0,1)\).
Given its class, its acquisition value \(\theta\) has a distribution \(G_j\) on \([0,\infty)\) with a finite
mean; atoms are allowed. Write \(G_j(p^-)\) for the probability that a class-\(j\) value falls below the
reserve \(p\), and \(D(y)=G_L(y)-G_H(y)\).

% STRUCTURE: Assumption 1 (class ranking) -> prose with the display eq:classes; label ass:classes retired
Two restrictions rank the classes. The high class, or high-value class, dominates the low class in the sense of
first-order stochastic dominance, and both class supports lie above the reserve, so that every value meets it:
\begin{equation}\label{eq:classes}
G_H(y)\le G_L(y)\quad\text{for every }y,\qquad G_H(p^-)=G_L(p^-)=0.
\end{equation}
I call these two restrictions the class ranking.

Under the class ranking~\eqref{eq:classes}, competition widens the class difference in target proceeds under every
incumbent distribution (Proposition~\ref{prop:returns}). Appendix~\ref{app:payoffs} characterizes the weakest
restriction of this kind. When neither class has a value atom at the reserve, the spread rises under every
strengthening exactly when two conditions hold. The high class dominates above the reserve, and the two classes
are equally likely to fall below it, which I call equal admissibility. Section~\ref{sec:reserves} studies a
reserve that excludes the low class, where equal admissibility fails.

The structure is the selective-entry structure of the auction literature. The class indexes the distribution of
the challenger's value, preparation reveals the value, and the stock price is the pre-entry signal.

The benchmark is the point-mass case: \(G_H\) puts all its mass on \(h\), \(G_L\) puts all its mass on
\(\ell\), and \(\pi=\reg{prim_pi}\). In the benchmark the class and the value coincide. Values, the reserve, and the
uniform incumbent's strength satisfy
\begin{equation}\label{eq:order}
0<p<\ell<r<h.
\end{equation}

The general statements need no normalization of money. A common positive rescale of every value, the reserve,
and every preparation, trading and research cost rescales prices and payoffs by the same factor. Orders,
beliefs, entry, and ownership do not change. The benchmark measures money in units of the low value, so
\(\ell=\reg{prim_ell}\) there.

The challenger initially knows neither its class nor its value. After it observes the stock price, it
privately learns its preparation cost \(C\in\{0,c_H\}\), with \(\Pr(C=0)=\rho\), \(0<\rho<1\), and
\(c_H>0\). The low cost is zero: a zero-cost challenger is a ready buyer whose diligence costs nothing at the
margin. A payment of \(C\) reveals \(\theta\) and permits a bid. A challenger that does not pay stays outside
the sale.

The pair \((j,\theta)\), the incumbent's value, the preparation cost, and the draws of the order slot in
Section~\ref{sec:slot} are mutually independent. The cost therefore carries no information about the class or
the value. Conditional on the class, the challenger's value is independent of the order.

Before trade, the seller publicly commits to a cash second-price auction with reserve \(p>0\). The highest
admissible bidder acquires the target. It pays the larger of the reserve and the highest competing admissible
bid. A bid equal to the reserve is admissible, and there is no sale without an admissible bid. Each prepared
bidder whose value meets the reserve bids its value. This bid is weakly dominant (Appendix~\ref{app:payoffs}).

A prepared challenger whose value is below the reserve does not bid. Without a challenger, the incumbent buys
at the reserve when its value meets it. Expected no-entry proceeds are therefore \(t_0=p[1-F(p)]\): the board
collects its announced walk-away price from a lone prepared bidder whose value meets it. Because \(F\) is
continuous, acquisition-value ties have probability zero.

\subsection{Trading: the order slot}\label{sec:slot}

An investor observes the challenger's class \(j\). It is an outside sector specialist, such as a
sector-focused long-short fund. It holds no net position in the target before trade, and a holder with a full
hedge counts as an outsider in this sense. Its information concerns the challenger's fit with the target. It
forms that view from its own research on the target's customers, technology, and demand. Neither bidder nor the
target supplies this information.

Trade takes place in one order slot:
\begin{enumerate}[label=(\arabic*),leftmargin=2.5em]
\item One public order opportunity comes before the challenger's decision.
\item With probability \(\lambda\in(0,1)\), the informed investor fills the slot. The investor sees the class
\(j\) and chooses an order \(q\in\{-1,0,1\}\). A nonzero order costs \(k>0\).
\item With probability \(1-\lambda\), a liquidity trader fills the slot. It sells, abstains, or buys one unit
with probabilities \(\beta\), \(1-2\beta\), and \(\beta\), where \(0<\beta\le1/3\). Its motive is exogenous.
\item No one observes who filled the slot. The identity draw is independent of the class, the costs, and the
incumbent's value.
\item Competitive market makers see the order \(X\in\{-1,0,1\}\) and set the price
\(P=\E[V_T\mid X]\). The challenger sees the price and then decides whether to prepare.
\end{enumerate}
% STRUCTURE: the symbol eta for the liquidity order distribution -> Appendix A opening, with Definition A.1

If the investor fills the slot with order \(q\), its profit is
\begin{equation}\label{eq:profit}
q\,[V_T-P(q)]-k|q|,
\end{equation}
where \(V_T\) is the terminal payoff of a target share. The investor's order is the whole order that market
makers see, so it trades at the price \(P(q)\). The unit is a normalized small trading unit, not ownership of
the target. The cost \(k\) is a separate friction, such as an execution cost or a cost to carry the position. It is
distinct from adverse-selection price impact, which the competitive price already generates. The investor's
order after bad news is a short sale, and \(k\) includes the cost to borrow shares.

Liquidity traders fill the slot when the investor does not. A current shareholder sells one unit for cash, an
outside investor buys one unit, or the opportunity passes with no order. Each acts for reasons apart from the
sale. Competitive market makers absorb the order at the price \(P(X)\). Each share held at the sale,
market-maker inventory included, receives \(V_T\).

Two numbers summarize the slot. The probability of a liquidity order in a given direction is
\(b=(1-\lambda)\beta\). The ratio \(u=\lambda/b\) is the informed presence mass relative to the liquidity mass in
one direction. The number \(1+u\) is the largest likelihood ratio that one order can carry. Section~\ref{sec:inference} shows that,
under every investor strategy, the posterior after every order lies in \([m(\pi),M(\pi)]\), where
\begin{equation}\label{eq:bounds}
M(\pi)=\frac{\pi(1+u)}{\pi(1+u)+1-\pi},\qquad
m(\pi)=\frac{\pi}{\pi+(1-\pi)(1+u)}.
\end{equation}

The bounds follow from one property of the slot. Each class reaches each order through a liquidity trader, so
no order is conclusive about the class. The investor therefore keeps a residual advantage after every order.

The slot uses the alternative-arrival structure of \citet{DowGorton1997}: an informed order and a liquidity
order of the same size are alternative arrivals in one opportunity. In \citet{AdmatiPfleiderer2009}, a
liquidity shock forces the informed holder's own sale, so a sale occurs on path for both motives. The unit order
menu, the presence probability, and the symmetric unit cost follow \citet{EdmansGoldsteinJiang2015}.%
% STRUCTURE: Remark 3 (continuous order sizes) and Section 6.4 -> this footnote; labels rem:sizes and sec:sizes retired
\footnote{The unit menu \(\{-1,0,1\}\) keeps every order on path, so no off-path belief enters the analysis. A menu
of continuous order sizes needs a convention for the sizes that no class uses in equilibrium, as in
\citet{EdmansGoldsteinJiang2015}. Another route keeps every size on path through additive liquidity flow. With
Laplace noise of scale \(s\) and full unit orders, the largest likelihood ratio of a favorable flow is
\(\eu^{2/s}\), so the top posterior equals the slot's \(M(\pi)\) when \(1+u=\eu^{2/s}\). The two set-ups then
share the posterior bounds. They differ in the probability with which each class reaches the top posterior.}
Risk-neutral market makers with zero expected profit price each order, as in \citet{GlostenMilgrom1985} and
\citet{Kyle1985}. Here the terminal payoff depends on the price through preparation.

Market makers observe the order \(X\) and nothing else. The challenger observes the price \(P\) and its own cost
\(C\), but not \(X\), its class, its value or the incumbent's value. The investor observes the class \(j\) but not
the value \(\theta\), and the incumbent knows its own value. The sequence is as follows:
\begin{enumerate}[label=(\arabic*),leftmargin=2.5em]
\item The seller announces the sale rule.
\item The investor learns the class. The slot draw selects the investor with probability \(\lambda\) and a
liquidity trader otherwise, and the selected trader places its order.
\item Market makers set the price from the order.
\item The challenger observes the price and its cost and decides whether to prepare.
\item Prepared bidders learn their values and submit their bids.
\item Ownership and financial payoffs are realized.
\end{enumerate}
Section~\ref{sec:research} adds a research stage before step (2).

\subsection{Equilibrium}\label{sec:equilibrium}

For a known challenger value \(\theta\), write \(t(\theta)\) for expected target proceeds after the challenger
prepares and \(g(\theta)\) for the challenger's expected gross acquisition profit. As before, \(t_0\) is expected
target proceeds without entry. The class objects are expectations over the class distribution,
\[
t_j=\E_j[t(\theta)],\qquad g_j=\E_j[g(\theta)],\qquad j\in\{H,L\},
\]
and the target-payoff spread is \(\DT=t_H-t_L\). Bids are truthful, so these objects depend on the incumbent's
distribution and the reserve but not on any belief.

At a belief \(\mu=\Pr(j=H)\), the challenger's expected gross profit is
\[
B(\mu)=g_L+\mu\,(g_H-g_L),
\]
because its value depends on the price only through its class. When the incumbent's distribution matters, I
write \(g_j(F)\), \(\DT(F)\), and \(B_F\). For the uniform incumbent on \([0,r]\), I write \(g_j(r)\),
\(\DT(r)\), and \(B_r\).

% STRUCTURE: Definition 1 (equilibrium tuple, conditions (E1)-(E6)) -> Appendix A opening, Definition A.1
An equilibrium specifies, for each class, a distribution over the three orders, the market makers' posterior
\(\mu(x)=\Pr(j=H\mid X=x)\) and price \(P(x)\) after each order \(x\), the challenger's belief at each price, its
preparation choice at each price and cost, and truthful bids. The investor can mix, and wrong-signed orders are
allowed. Every order has positive probability in both classes, because a liquidity trader can place it, so Bayes'
rule defines the posterior after every order and no off-path belief enters. Prices are competitive, a high-cost
challenger prepares when gross profit at its belief exceeds \(c_H\) and stays out when it falls short, and a
zero-cost challenger, whose gross profit is an expected auction rent and never negative, prepares after every
price. The investor's order maximizes
\[
U_j(q)=q\,\bigl[t_0+e(q)(t_j-t_0)-P(q)\bigr]-k|q|
\]
at the candidate's prices and entry, where \(e(q)\) is entry after the order \(q\), the probability that the
challenger prepares after the price \(P(q)\). Definition~\ref{def:equilibrium} in Appendix~\ref{app:proofs} states
the six conditions in full.

The investor's objective uses the expected terminal payoff of a target share, given class \(j\) and order \(q\).
Entry depends only on the price and the cost, and the cost is independent of the class and the order. That payoff
is therefore \(t_0+e(q)(t_j-t_0)\). The same independence and Bayes' rule turn the competitive price into
\begin{equation}\label{eq:price}
P(x)=t_0+e(x)\,\bigl[t_L-t_0+\DT\,\mu(x)\bigr].
\end{equation}
The ex ante payoff of the investor multiplies \(U_j\) by \(\lambda\). This factor leaves the best order
unchanged.

Two indifferences remain: an investor can be indifferent among orders, and a high-cost challenger can be
indifferent about preparation. The equilibrium concept lets an indifferent investor mix over its best orders in
any proportion. This is the standard best-response condition, and I call it rule IM. An indifferent high-cost
challenger prepares, and I call this rule C1. The results use IM and C1.
% STRUCTURE: the list of the alternatives I0 and CM, and the sentence on the remarks -> Appendix A.4, paragraph "Other tie rules"
Appendix~\ref{app:coexistence} states two alternatives, rule I0, under which an investor indifferent between
abstention and a nonzero order abstains, and rule CM, under which an indifferent high-cost challenger may prepare
with any probability; Table~\ref{tab:a4-ties} records which conclusions on the benchmark path depend on them.

Two comparisons recur. A unilateral deviation by the investor changes only its own order. The price function
\(P\) and the entry rule stay at their values in the candidate equilibrium. The deviation payoff is
therefore \(U_j(q)\) read at the candidate's \(P(q)\) and \(e(q)\). A comparison across economies, such as
across incumbent strengths, solves for the equilibrium again in each economy. Proofs use the first comparison to
test a candidate equilibrium and the second to compare outcomes.

Two outcomes summarize an equilibrium. \emph{Entry} \(\Ent\) is the probability that the challenger prepares.
\emph{High-class ownership} \(\Own_H\) is the probability that a challenger of class \(H\) acquires the target.
Write \(W_H(F)=\E_H\bigl[F(\theta)\1\{\theta\ge p\}\bigr]\) for the probability that a prepared challenger of
class \(H\) wins against the incumbent distribution \(F\).
